% ============================================================
% Command-Line Options
% ============================================================

\begin{longtable}{p{0.34\textwidth} p{0.58\textwidth}}
    \caption{Command-line options available when starting LuRaCs.}
    \label{tab:cli-options} \\

    \hline
    \textbf{Option} &
    \textbf{Description} \\
    \hline
    \endfirsthead

    \multicolumn{2}{c}%
    {\tablename\ \thetable{} -- continued from previous page} \\
    \hline
    \textbf{Option} &
    \textbf{Description} \\
    \hline
    \endhead

    \hline
    \multicolumn{2}{r}{Continued on next page} \\
    \endfoot

    \hline
    \endlastfoot

    % ------------------------------------------------------------
    % General
    % ------------------------------------------------------------

    \multicolumn{2}{l}{\textbf{General}} \\

    \texttt{--debug} &
        Launch the application in debug mode. A mock device is added
        automatically for testing and development. \\

    \texttt{--headless} &
        Run the application without the graphical user interface,
        using terminal mode instead. \\

    \texttt{--clear\_local\_cache} &
        Remove all locally stored files and exit the application. \\

    % ------------------------------------------------------------
    % Data Import
    % ------------------------------------------------------------

    \multicolumn{2}{l}{\textbf{Data Import}} \\

    \texttt{--import\_spectrum <file> [file ...]} &
        Load one or more spectrum files when starting the application.
        Each specified file must exist and be supported by the
        application's file importer. \\

    \texttt{--import\_rois <file>} &
        Load regions of interest from the specified XML file. \\

    \texttt{-roi <lower> <upper>} &
        Add a region of interest with the specified lower and upper
        energy bounds in keV. The option can be specified multiple
        times to add several ROIs. \\

    \texttt{--nuclides <nuclide> [nuclide ...]} &
        Preset one or more nuclides in the isotopics interface.
        This option is applied when the graphical user interface is
        available. \\

    % ------------------------------------------------------------
    % Device Connections
    % ------------------------------------------------------------

    \multicolumn{2}{l}{\textbf{Device Connections}} \\

    \texttt{-ble <device> [device ...]} &
        Attempt to establish Bluetooth Low Energy connections using
        the specified device names. \\

    \texttt{--bluetooth <device> [device ...]} &
        Long-form alias for \texttt{-ble}. Attempts Bluetooth Low
        Energy connections using the specified device names. \\

    \texttt{-usb <device> [device ...]} &
        Attempt to establish USB connections using the specified
        device names. Matching is performed against detected USB
        device product names. \\

    \texttt{-nw <ip> [ip ...]} &
        Attempt to connect to network devices using their IP addresses. \\

    \texttt{--network <ip> [ip ...]} &
        Long-form alias for \texttt{-nw}. Connects to network devices
        using their IP addresses. \\

    \texttt{--serial\_gps <interface>} &
        Connect a GPS receiver using a USB serial interface. The
        specified value identifies the serial interface used by
        the GPS receiver. \\

    % ------------------------------------------------------------
    % Maps
    % ------------------------------------------------------------

    \multicolumn{2}{l}{\textbf{Maps}} \\

    \texttt{--map\_url <URL>} &
        Load an online map from the specified URL at application startup. \\

    \texttt{--map\_file <path>} &
        Load a local offline map from the specified \texttt{.pmtiles}
        file at application startup. \\

    \texttt{--map\_url} and \texttt{--map\_file} &
        Only one map source may be specified at startup. If both
        options are provided, no map is loaded and a warning is
        written to the application log. \\

\end{longtable}
